\chapter{Robot dan Kendaraan Otonom}
\label{ch:robot}

Musim panas 2025 membawa saya ke dua kuliah tentang mesin yang bergerak. Introduction to robotic
systems membahas mekanika lengan robot: kerangka koordinat, kinematika maju dan balik, Jacobian,
persamaan gerak dari prinsip Lagrange, visual servoing, dan perencanaan jalur. Introduction to
autonomous vehicles membahas mobil yang mengemudi sendiri: lokalisasi dengan pendekatan Markov, ICP,
dan NDT, sensor lidar, deteksi objek tiga dimensi, Kalman filter, dan pelacakan banyak objek, dengan
proyek lokalisasi dan proyek persepsi.

Robot adalah tempat semua yang sudah kita bahas bertemu. Ia harus melihat (\cref{ch:cv}), menaksir
keadaannya di tengah ketidakpastian (\cref{ch:prob}), mengendalikan geraknya (\cref{ch:kendali}),
dan memutuskan apa yang dilakukan (\cref{ch:rl,ch:klasik}). Bab ini melihat dua jenis robot:
lengan industri yang bekerja di tempat, dan kendaraan yang bergerak di dunia terbuka.

\section{Koordinat dan kinematika}

Robot hidup di antara banyak kerangka koordinat: dunia, dasar robot, setiap sendi, ujung lengan,
kamera. Posisi dan orientasi sebuah kerangka relatif terhadap kerangka lain dinyatakan dengan
matriks rotasi dan vektor translasi, yang digabung menjadi matriks transformasi homogen $4\times4$,
sehingga transformasi berurutan cukup dikalikan. Konvensi Denavit--Hartenberg \citep{denavit1955}
memberi resep standar untuk menempatkan kerangka di setiap sendi, sehingga setiap sambungan hanya
butuh empat parameter.

\istilah{Forward kinematics} menghitung posisi ujung lengan dari sudut-sudut sendi. Untuk lengan planar
dua ruas dengan panjang $l_1$ dan $l_2$ serta sudut $q_1$ dan $q_2$,
\begin{equation}
x=l_1\cos q_1+l_2\cos(q_1+q_2),\qquad y=l_1\sin q_1+l_2\sin(q_1+q_2).
\end{equation}
Dengan $l_1=l_2=1$, $q_1=0$, dan $q_2=90^\circ$, ujung lengan berada di $(1,1)$. \istilah{Inverse kinematics}
berjalan ke arah sebaliknya: dari posisi yang diinginkan, cari sudut sendinya. Untuk lengan dua ruas ada
dua jawaban, siku di atas atau siku di bawah, dan untuk posisi di luar jangkauan tidak ada jawaban sama
sekali. Lengan dengan lebih banyak sendi dari yang dibutuhkan, yang disebut redundan, punya tak hingga
jawaban, dan kebebasan itu bisa dipakai untuk menghindari rintangan.

Jacobian menghubungkan kecepatan sendi dengan kecepatan ujung lengan, $\dot{\bm p}=\bm J(\bm q)\dot{\bm q}$.
Untuk lengan dua ruas tadi, determinannya $l_1l_2\sin q_2$. Ketika lengan terbentang lurus atau
terlipat penuh, $\sin q_2=0$ dan Jacobian kehilangan pangkat. Di konfigurasi \istilah{singular} seperti
ini, ada arah gerak yang tidak bisa dicapai sama sekali, dan kecepatan sendi yang dibutuhkan untuk
bergerak di dekatnya membengkak.

\section{Dinamika}

Kinematika tidak peduli pada massa dan gaya. Untuk mengendalikan lengan yang sebenarnya, kita perlu
tahu torsi apa yang menghasilkan gerak tertentu. Kuliah menurunkannya dengan prinsip Hamilton dan
persamaan Lagrange. Tuliskan energi kinetik $T$ dan energi potensial $V$ sebagai fungsi koordinat sendi,
bentuk fungsi Lagrange $L=T-V$, lalu terapkan
\begin{equation}
\frac{\dd}{\dd t}\frac{\partial L}{\partial\dot q_i}-\frac{\partial L}{\partial q_i}=\tau_i.
\end{equation}
Hasilnya selalu berbentuk $\bm M(\bm q)\ddot{\bm q}+\bm C(\bm q,\dot{\bm q})\dot{\bm q}+\bm g(\bm q)=\bm\tau$,
dengan matriks massa yang bergantung pada konfigurasi, suku Coriolis dan sentrifugal, dan gravitasi.
Struktur ini dipakai langsung untuk kendali: kalau persamaannya diketahui, pengendali bisa menghitung
torsi yang membatalkan nonlinearitas dan menyisakan sistem linear yang mudah dikendalikan. Deskripsi
robot untuk perangkat lunak, seperti format URDF, menyimpan semua massa, inersia, dan geometri yang
dibutuhkan untuk perhitungan ini \citep{lynch2017}.

\section{Melihat dan merencanakan}

Visual servoing mengendalikan robot dengan umpan balik dari kamera. Model kamera dan kalibrasinya sama
dengan di \cref{ch:cv}. Pada visual servoing berbasis posisi, pose objek diperkirakan dari gambar, lalu
galat pose diumpankan ke pengendali.

Untuk bergerak tanpa menabrak, kuliah membahas medan potensial buatan. Tujuan memberi potensial yang
menarik, rintangan memberi potensial yang menolak, dan robot mengikuti gradien negatif jumlahnya, seperti
bola yang menggelinding di lanskap. Cara ini elegan dan murah, tetapi punya masalah yang sudah kita
kenal dari optimasi: minimum lokal. Robot bisa terjebak di lekukan antara dua rintangan, di mana tarikan
dan tolakan seimbang. Perencana berbasis graf seperti A* (\cref{ch:klasik}) atau perencana berbasis sampel
menghindarinya dengan mencari secara global. Gerak sepanjang jalur kemudian diberi profil kecepatan
yang halus, misalnya trapesium dengan fase percepatan, kecepatan tetap, dan perlambatan, supaya motor
tidak dihentak.

\section{Di mana aku berada: lokalisasi}

Mobil otonom harus tahu posisinya dengan ketelitian beberapa sentimeter, jauh lebih teliti dari GPS biasa.
Kuliah membahas tiga pendekatan. \istilah{Markov localization} adalah filter Bayes dari \cref{ch:prob} di atas
grid posisi \citep{thrun2005}. Keyakinan tentang posisi disebar menurut model gerak, karena roda selip
dan odometri tidak sempurna, lalu dipertajam dengan membandingkan pengukuran sensor dengan peta.
Bayangkan berjalan di koridor gelap dengan tangan meraba dinding: setiap langkah menambah keraguan,
setiap pintu yang teraba mengurangi keraguan.

Dua pendekatan lain bekerja dengan mencocokkan awan titik lidar dengan peta. \istilah{Iterative closest
point} \citep{besl1992} bergantian antara memasangkan setiap titik pindaian dengan titik peta terdekat,
lalu mencari rotasi dan translasi kaku yang meminimalkan jarak kuadrat pasangan-pasangan itu. Langkah
kedua punya solusi tertutup melalui dekomposisi nilai singular. ICP cepat konvergen kalau tebakan
awalnya baik, dan bisa tersesat kalau tidak. \istilah{Normal distributions transform} \citep{biber2003}
membagi peta menjadi sel-sel dan mewakili titik-titik di setiap sel dengan sebuah distribusi Gauss.
Pindaian baru dicocokkan dengan memaksimalkan likelihood titik-titiknya di bawah campuran Gauss itu,
yang lebih mulus dan lebih tahan terhadap tebakan awal yang kurang tepat.

\section{Melihat jalan: persepsi dan pelacakan}

Lidar memancarkan pulsa laser dan mengukur waktu tempuhnya, sehingga menghasilkan awan titik tiga
dimensi dengan jarak yang akurat. Kamera memberi warna dan tekstur, dan radar (\cref{ch:sinyal})
memberi kecepatan langsung dan tetap bekerja dalam hujan dan kabut. Fusi sensor menggabungkan
kelebihan masing-masing.

Mendeteksi objek dalam awan titik punya kesulitan khas: titik-titiknya tidak tersusun dalam grid dan
urutannya tidak bermakna. PointNet \citep{qi2017} menanganinya dengan menerapkan MLP yang sama pada
setiap titik lalu menggabungkan hasilnya dengan pooling maksimum, sebuah fungsi yang invarian terhadap
permutasi persis seperti DeepSets di \cref{ch:gdl}. PointPillars \citep{lang2019} mengelompokkan titik
ke dalam pilar-pilar vertikal di atas grid tanah, mengubahnya menjadi gambar pandangan atas, lalu
memakai jaringan konvolusi dua dimensi yang cepat.

Deteksi di setiap bingkai belum cukup. Kita perlu tahu bahwa mobil di bingkai ini adalah mobil yang
sama dengan di bingkai sebelumnya, dan ke mana ia bergerak. Pelacakan banyak objek bekerja dalam
siklus. Untuk setiap jejak yang sudah ada, Kalman filter atau versi diperluasnya meramalkan posisi di
bingkai berikutnya. Deteksi baru lalu dipasangkan dengan jejak-jejak itu. Pemasangan ini adalah masalah
penugasan: diberikan matriks biaya berisi jarak antara setiap ramalan dan setiap deteksi, cari pasangan
satu-satu dengan biaya total terkecil. Algoritme Hungaria \citep{kuhn1955} menyelesaikannya secara eksak.
Deteksi yang tidak punya pasangan bisa menjadi jejak baru, dan jejak yang lama tidak mendapat deteksi
dihapus. Setiap jejak yang terpasang diperbarui dengan langkah koreksi Kalman.

\section{Di mana machine learning masuk}

Robotika klasik dibangun dari model yang diturunkan: kinematika, dinamika, model sensor. Machine learning
masuk di tempat model sulit ditulis. Persepsi hampir seluruhnya dikuasai deep learning. Dinamika yang
sulit dimodelkan, seperti gesekan, kabel, atau benda lunak, bisa dipelajari dari data dan dipakai
sebagai koreksi atas model fisika. Policy untuk manipulasi yang rumit dilatih dengan reinforcement
learning di simulator lalu dipindahkan ke robot nyata, atau dipelajari dari demonstrasi manusia
(\cref{ch:rl}). Model fondasi yang memadukan penglihatan, bahasa, dan aksi mulai dipakai untuk memberi
perintah robot dengan kalimat biasa (\cref{ch:terbaru}). Tetapi untuk bagian yang menyangkut keselamatan,
seperti lokalisasi, pelacakan, dan kendali tingkat rendah, metode klasik dengan jaminan yang bisa
dianalisis tetap menjadi tulang punggung.

\sumberbab{Bab ini mengikuti kuliah Introduction to Robotic Systems (transformasi kerangka, kecepatan
benda kaku, Denavit--Hartenberg, kinematika maju dan balik, Jacobian dan singularitas, persamaan
Lagrange, visual servoing, medan potensial, profil trayektori) dan Introduction to Autonomous Vehicles
(Markov localization, ICP, NDT, lidar, deteksi objek tiga dimensi, Kalman dan extended Kalman filter,
pelacakan banyak objek). Buku rujukan: \citet{lynch2017} dan \citet{thrun2005}.}
